% 第6章 図6-2: 製作費が交絡になる場合 (左) と通り道になる場合 (右)
% 箱と数字は同じで, 違うのは矢印の向きだけ
\begin{tikzpicture}[
    font=\small,
    box/.style={draw, rounded corners=2pt, minimum width=2.1cm, minimum height=0.75cm, fill=black!5},
    arr/.style={-stealth, thick, shorten >=1pt},
  ]
  % 左: 製作費が原作の有無と視聴回数の両方に先立つ
  \node[box] (L-cost) at (0,1.6) {製作費};
  \node[box] (L-src)  at (-1.9,0) {原作の有無};
  \node[box] (L-view) at (1.9,0) {視聴回数};
  \draw[arr] (L-cost) -- (L-src);
  \draw[arr] (L-cost) -- (L-view);
  \draw[arr] (L-src) -- (L-view);
  \node[font=\footnotesize] at (0,-0.9) {交絡になる場合};
  % 右: 原作から製作費を通って視聴回数へ
  \begin{scope}[xshift=7.2cm]
    \node[box] (R-cost) at (0,1.6) {製作費};
    \node[box] (R-src)  at (-1.9,0) {原作の有無};
    \node[box] (R-view) at (1.9,0) {視聴回数};
    \draw[arr] (R-src) -- (R-cost);
    \draw[arr] (R-cost) -- (R-view);
    \draw[arr] (R-src) -- (R-view);
    \node[font=\footnotesize] at (0,-0.9) {通り道になる場合};
  \end{scope}
\end{tikzpicture}
